%TCIDATA{Version=5.50.0.2953}
%TCIDATA{LaTeXparent=0,0,main.tex}
                      

\pretolerance=500000 \tolerance=500000

\appendix{}

\chapter{Informaci\'{o}n complementaria}

Los \textbf{ap\'{e}ndices} se colocan al final de un documento acad\'{e}%
mico, generalmente no se incluyen en el recuento total de palabras y
proporcionan material de apoyo a la investigaci\'{o}n; algunos ejemplos se
presentan a continuaci\'{o}n:

\begin{enumerate}
\item Cartas relevantes a participantes y organizaciones; por ejemplo, actas
de comit\'{e}s de \'{e}tica o declaraciones de consentimiento informado.

\item Informes de antecedentes.

\item Datos brutos.

\item C\'{a}lculos detallados.

\item Experimentos complementarios.

\item Diferentes tipos de datos se presentan en ap\'{e}ndices separados.
\end{enumerate}

Cada ap\'{e}ndice debe tener un t\'{\i}tulo, estar etiquetado con un n\'{u}%
mero o letra y se debe hacer referencia a este en el cuerpo del documento.

\bigskip

Un buen \textbf{t\'{\i}tulo} debe ser claro y breve, con la menor cantidad
de palabras posibles, pero describir con certeza el contenido significativo
del trabajo, debe transmitir su alcance y atraer a los potenciales lectores;
casi nunca deben contener abreviaturas, f\'{o}rmulas, nombres patentados (en
lugar de gen\'{e}ricos) o jerga (lenguaje informal); el t\'{\i}tulo de un
documento acad\'{e}mico es una etiqueta, normalmente no es una oraci\'{o}n,
es m\'{a}s simple o, al menos, m\'{a}s corta, pero el orden de las palabras
es m\'{a}s importante; por lo tanto, la sintaxis es otro elemento que se
debe considerar, algunos errores gramaticales en los t\'{\i}tulos son
resultado del orden incorrecto de las palabras. Por lo general, tanto el t%
\'{\i}tulo como el resumen se desarrollan al final, ya que todo el cuerpo
del documento est\'{a} escrito, sin importar que estos se encuentren al
inicio del mismo. El t\'{\i}tulo suele ser la primera introducci\'{o}n que
los lectores (y revisores) tienen del trabajo, por lo tanto, se debe
seleccionar uno que capte la atenci\'{o}n, describa con precisi\'{o}n el
contenido y haga que la comunidad cient\'{\i}fica quiera leerlo.

\bigskip

La lista de \textbf{referencias} contiene todos los recursos que se han
citado en el desarrollo del trabajo, estas fuentes pueden incluir libros, art%
\'{\i}culos, sitios web, documentos t\'{e}cnicos y otros materiales que se
utilizaron para recopilar informaci\'{o}n. A continuaci\'{o}n, se presentan
algunas pautas para crear una lista de referencias:

\begin{enumerate}
\item Formatear las citas: Las referencias deben presentarse de acuerdo con
un estilo o formato de cita espec\'{\i}fico; por ejemplo, APA, MLA o IEEE,
entre otros.

\item Incluir toda la informaci\'{o}n necesaria: Cada fuente debe de tener
el nombre de los autores, el t\'{\i}tulo de la obra, la fecha de publicaci%
\'{o}n, la editorial, nombre de la revista y cualquier otra informaci\'{o}n
relevante como v\'{\i}nculos, doi (del ingl\'{e}s \textit{digital object
identifier}), volumen y n\'{u}mero, p\'{a}ginas, n\'{u}mero de edici\'{o}n,
etc.
\end{enumerate}

Es importante recordar que una bibliograf\'{\i}a contiene todas las fuentes
consultadas para la investigaci\'{o}n, mientras que la lista de referencias
se limita a las fuentes espec\'{\i}ficamente citadas dentro del texto.

\bigskip

Para informaci\'{o}n adicional sobre la escritura de documentos cient\'{\i}%
ficos, informes t\'{e}cnicos y art\'{\i}culos de investigaci\'{o}n, se
sugiere leer las siguientes referencias:

\begin{enumerate}
\item El libro \textquotedblleft \textit{How to write and publish a
scientific paper}\textquotedblright\ de Gastel y Day \cite{Gastel(2022)}.

\item El art\'{\i}culo \textquotedblleft \textit{Writing a scientific
article: A step-by-step guide for beginners}\textquotedblright\ de Ecarnot 
\textit{et al}. \cite{Ecarnot(2015)}.

\item El art\'{\i}culo de p\'{a}gina web \textquotedblleft \textit{Redacci%
\'{o}n de un manuscrito de revista}\textquotedblright\ de la editorial
Springer \cite{springer}.

\item El art\'{\i}culo de p\'{a}gina web \textquotedblleft \textit{Research
reports}\textquotedblright\ de la Universidad de Melbourne \cite{reports}.
\end{enumerate}

Otros documentos que se pueden utilizar como referencias en el desarrollo y
escritura de libros, art\'{\i}culos de investigaci\'{o}n, art\'{\i}culos de
revisi\'{o}n, tesis de licenciatura y posgrado, disertaciones de ingenier%
\'{\i}a, informes t\'{e}cnicos, y memorias de conferencias nacionales e
internacionales son los siguientes:

\begin{enumerate}
\item Tesis de Doctorado \cite{ValleScD}.

\item Tesis de Maestr\'{\i}a \cite{ValleMC}.

\item Tesis de Licenciatura \cite{ValleENG}.

\item Residencias Profesionales \cite{ValleRES}.

\item Cap\'{\i}tulos de Libros \cite{Byers(2009)}.

\item Conferencias en Simposios o Congresos \cite{BEER(2020)}.

\item Art\'{\i}culos en Memorias de Congresos \cite{Korobeinikov(2017)}.

\item \textit{Syllabus} o Manuales de Pr\'{a}cticas \cite{biomath}.
\end{enumerate}
